% ==============================================================================
% Section 5: Conclusion & Mathematical Reference Summary
% ==============================================================================

\clearpage
\section[Conclusion \& Formula Reference]{Conclusion and Formula Reference Summary}

\subsection{Synthesis of Findings and Strategic Investment Implications}
The quantitative empirical analysis of ten marquee Indian FMCG corporations illustrates the profound dichotomy between static asset liquidation values and forward-looking market pricing:

\begin{enumerate}[label=\textbf{\arabic*.}, topsep=3pt, itemsep=2pt, parsep=0pt]
    \item \textbf{Uniform Overvaluation under the Liquidation Metric:} Every company examined trades at a substantial multiple above its net tangible asset worth, ranging from $4.5\times$ (ITC Ltd) to $42.7\times$ (Nestle India Ltd).
    \item \textbf{Dual Investor Perspectives:}
    \begin{itemize}[topsep=2pt, itemsep=2pt, parsep=0pt]
        \item \textbf{The Classical Value Investor Perspective:} Under strict Benjamin Graham criteria, none of the ten companies provide an adequate margin of safety based on liquidation assets alone. However, within the peer group, \textbf{ITC Ltd} (\inr 74,257 Cr net assets, $4.5\times$ ratio) and \textbf{Tata Consumer Products} (\inr 20,195 Cr net assets, $4.9\times$ ratio) offer superior downside balance sheet defense.
        \item \textbf{The Franchise / Quality Investor Perspective:} For growth and franchise investors, the market premium reflects durable competitive advantages (moats), high cash-flow conversion, dependable dividend distributions, and superior Return on Capital Employed ($\operatorname{ROCE}$). Companies like \textbf{Nestle India} and \textbf{Britannia Industries} trade at $33\times - 43\times$ tangible book value precisely because their enterprise value is propelled by high-margin product consumption rather than factory equipment.
    \end{itemize}
    \item \textbf{Valuation Synthesis:} While the Net Asset Value method fails as an independent estimator of fair market value for asset-light corporations, it performs a vital risk-management function by demarcating the \textbf{hard liquidation floor} of the firm in catastrophic downturns.
\end{enumerate}

\subsection{Master Reference Table: Valuation Equations}
Table~\ref{tab:formula_summary} provides a consolidated reference of all mathematical formulas implemented throughout this equity valuation project.

\begin{table}[H]
\centering
\small
\renewcommand{\arraystretch}{1.18}
\begin{tabularx}{\textwidth}{l p{5.5cm} X}
\toprule
\textbf{Valuation Metric} & \textbf{Mathematical Formulation} & \textbf{Operational Financial Purpose} \\
\midrule
\textbf{Average Profit} & $\overline{\Pi} = \dfrac{1}{5}\sum_{t=1}^{5} \Pi_t$ & Normalizes historical profitability over a 5-year cycle. \\
\textbf{Goodwill} & $\text{Goodwill} = \overline{\Pi} \times N_{\text{purchase}} = 3\,\overline{\Pi}$ & Quantifies capitalized super-profit earning power. \\
\textbf{Tangible Assets} & $\text{Tangible Assets} = \text{Total Assets} - \text{Goodwill}$ & Eliminates unamortized/subjective intangibles. \\
\textbf{Net Assets ($\NAV$)} & $\text{Net Assets} = \text{Tangible Assets} - \text{Total Liab.}$ & Measures residual tangible pool for common equity. \\
\textbf{Intrinsic Value ($\IV$)} & $\IV = \dfrac{\text{Net Assets Available}}{\text{Number of Equity Shares}}$ & Establishes per-share liquidation value floor. \\
\textbf{Market Price ($P_0$)} & $P_0 = \dfrac{\text{Market Capitalization}}{\text{Number of Equity Shares}}$ & Captures consensus valuation in secondary markets. \\
\textbf{Valuation Gap ($\Delta$)} & $\Delta = \IV - P_0$ & Quantifies absolute divergence from asset backing. \\
\textbf{Overvaluation Multiple} & $\mathcal{R} = \dfrac{P_0}{\IV} = \dfrac{\text{Market Cap}}{\text{Net Assets}}$ & Evaluates relative pricing premium over physical equity. \\
\bottomrule
\end{tabularx}
\caption{Consolidated Mathematical Formulations for Net Asset Value Equity Valuation.}
\label{tab:formula_summary}
\end{table}
